% !PRINTED 98
\weberpage{98}
%<<CONT
neation $\mathfrak{M}$ plays the role of the unit. If $S = \begin{pmatrix} a, & b \\ c, & \partial \end{pmatrix}$ is an arbitrary transformation, then
\[
\begin{pmatrix} a, & b \\ c, & \partial \end{pmatrix}\begin{pmatrix} \partial, & -b \\ -c, & a \end{pmatrix} = \begin{pmatrix} n, & 0 \\ 0, & n \end{pmatrix}
\]
is a multiplication. The two collineations
\[
C = \mathfrak{M}\begin{pmatrix} a, & b \\ c, & \partial \end{pmatrix}, \quad C^{-1} = \mathfrak{M}\begin{pmatrix} \partial, & -b \\ -c, & a \end{pmatrix}
\]
therefore give, on composition, $CC^{-1} = C^{-1}C = \mathfrak{M}$, and so are reciprocal to one another.

\begin{center}
\gesperrt{Accordingly the totality of the collineations forms a group.}
\end{center}

The transformations of degree 1 are called \gesperrt{linear transformations}. For these we denote the transformation numbers, for distinction, by the Greek letters $\alpha, \beta, \gamma, \delta$, so that
\[
L = \begin{pmatrix} \alpha, & \beta \\ \gamma, & \delta \end{pmatrix}, \quad \alpha\delta - \beta\gamma = 1
\]
means a linear transformation. The system $\mathfrak{L}$ of the linear transformations is a group contained in $\mathfrak{S}$; for if
\[
L = \begin{pmatrix} \alpha, & \beta \\ \gamma, & \delta \end{pmatrix}, \quad L' = \begin{pmatrix} \alpha', & \beta' \\ \gamma', & \delta' \end{pmatrix},
\]
then
\[
L'' = L'L = \begin{pmatrix} \alpha'\alpha+\beta'\gamma, & \alpha'\beta+\beta'\delta \\ \gamma'\alpha+\delta'\gamma, & \gamma'\beta+\delta'\delta \end{pmatrix} = \begin{pmatrix} \alpha'', & \beta'' \\ \gamma'', & \delta'' \end{pmatrix}
\]
is likewise linear, and, for given $L'$, $L''$, one can determine $L$ from the equations
\[
\begin{aligned}
\alpha'\alpha+\beta'\gamma &= \alpha'', & \alpha'\beta+\beta'\delta &= \beta'', \\
\gamma'\alpha+\delta'\gamma &= \gamma'', & \gamma'\beta+\delta'\delta &= \delta''
\end{aligned}
\]
uniquely. The unit of the group $\mathfrak{L}$ is the identical substitution $\begin{pmatrix} 1, & 0 \\ 0, & 1 \end{pmatrix}$, and every substitution $L$ has its reciprocal $L^{-1}$, as the composition
\[
LL^{-1} = \begin{pmatrix} \alpha, & \beta \\ \gamma, & \delta \end{pmatrix}\begin{pmatrix} \delta, & -\beta \\ -\gamma, & \alpha \end{pmatrix} = \begin{pmatrix} 1, & 0 \\ 0, & 1 \end{pmatrix}
\]
shows. The group $\mathfrak{L}$ is infinite and is \gesperrt{not commutative}.
